
The experiment was executed as a single run with seed = 42 on the Waterbirds training set (4,795 samples). No multi-seed replication was performed. All values below derive from the files \texttt{h-e1/experiment\textbackslash \_results.json} and \texttt{h-e1/04\textbackslash \_validation.md}.

\subsubsection{5.1 Proxy Signal Quality at T\_id = 5}

The three gate criteria at T\_id = 5 are:

\begin{table}[htbp]
\centering
\resizebox{\columnwidth}{!}{%
\begin{tabular}{llll}
\toprule
Metric & Value & Target & Status \\
\midrule
Minority/majority $\tilde{g}$ ratio & 8.805 & $\geq$ 3.0x & PASS \\
AUC (minority group prediction) & 0.914 & > 0.70 & PASS \\
balance\_deviation (top-25\% subset) & 0.379 & $\leq$ 0.10 & FAIL \\
\bottomrule
\end{tabular}
}
\end{table}

\textbf{Overall gate result: PARTIAL (2/3 criteria pass).} The experiment is classified as having a PARTIAL gate outcome; the gate type was MUST\_WORK, requiring all three criteria to be satisfied.

AUC = 0.914 indicates that $\tilde{g}$ alone discriminates minority group membership (G1 $\cup$ G2 vs. G0 $\cup$ G3) with high accuracy at T\_id = 5, without access to group labels. The ratio of 8.805 indicates that the mean normalized gradient norm of minority samples is approximately 8.8 times that of majority samples at this epoch.

Per-group $\tilde{g}$ values at T\_id = 5:

\begin{table}[htbp]
\centering
\resizebox{\columnwidth}{!}{%
\begin{tabular}{llll}
\toprule
Group & Description & $\tilde{g}$ mean & g\_raw mean \\
\midrule
G0 & Landbird + Land (majority) & 0.0221 & 0.553 \\
G1 & Landbird + Water (minority) & 0.3126 & 8.100 \\
G2 & Waterbird + Land (minority) & 0.4328 & 10.487 \\
G3 & Waterbird + Water (majority) & 0.0936 & 2.350 \\
\bottomrule
\end{tabular}
}
\end{table}

G2 --- the smallest minority group at 56 training samples --- shows the highest mean $\tilde{g}$ (0.4328), consistent with the prediction that samples in minority conditions with no available shortcut maintain the largest prediction residuals.

\begin{figure}[htbp]
\centering
\includegraphics[width=\columnwidth]{figures/gate_metrics.png}
\caption{Gate metrics at T\_id=5}
\label{fig:gate_metrics}
\end{figure}

\textit{Figure 1. Bar chart of the three gate criteria at T\_id = 5. Left: minority/majority $\tilde{g}$ ratio (8.8x vs. 3.0x target). Center: AUC (0.914 vs. 0.70 target). Right: balance\_deviation (0.379 vs. 0.10 target; see Section 5.2 for interpretation).}

\begin{figure}[htbp]
\centering
\includegraphics[width=\columnwidth]{figures/distribution_epoch5.png}
\caption{Distribution of $\tilde{g}$ per group at epoch 5}
\label{fig:distribution_epoch5}
\end{figure}

\textit{Figure 2. Distribution of normalized gradient norms per group at T\_id = 5. Minority groups G1 and G2 are shifted to higher $\tilde{g}$ values relative to majority groups G0 and G3.}

\subsubsection{5.2 The balance\_deviation Failure: Criterion Design Analysis}

The balance\_deviation criterion measures how closely the top-25\% high-$\tilde{g}$ subset approximates a class-uniform distribution. The measured value (0.379) exceeds the $\leq$ 0.10 threshold by a factor of 3.8.

Post-hoc analysis in the reflection report identifies this as a criterion design mismatch rather than a mechanism failure. Minority groups G1 and G2 constitute approximately 5\% of the 4,795 training samples (240 samples). Even if the top-25\% subset (1,199 samples) captured 100\% of minority samples, the minority fraction within the subset would be at most 240/1,199 $\approx$ 20\%, far below the 50\% required for class balance. The $\leq$ 0.10 balance\_deviation target is therefore unachievable for any minority-focused selection method on this dataset.

The distinction drawn in the reflection report is between class balance (the fraction of each class in the selected subset) and minority recall (the fraction of true minority samples captured in the subset). The DFR-style retraining procedure that motivates subset construction requires minority enrichment, not class balance. Figure 4 (balance heatmap) shows that the top-25\% subset is strongly minority-enriched relative to the overall training distribution. The AUC = 0.914 provides indirect evidence that minority recall in the top-25\% is high, though direct minority recall was not computed in H-E1; a revised experiment (h-e1-v2) with minority recall as the criterion was proposed but not yet executed.

\begin{figure}[htbp]
\centering
\includegraphics[width=\columnwidth]{figures/balance_heatmap.png}
\caption{Balance heatmap of top-25\% subset}
\label{fig:balance_heatmap}
\end{figure}

\textit{Figure 4. Group composition of the top-25\% high-$\tilde{g}$ subset compared to the full training set. Minority groups G1 and G2 are overrepresented in the selected subset relative to their training-set proportions.}

\subsubsection{5.3 Temporal Persistence of the Proxy Signal}

Gradient norms were collected at epochs {1, 3, 5, 10}. Training loss and per-epoch proxy metrics:

\begin{table*}[htbp]
\centering
\resizebox{\dimexpr 2\columnwidth+\columnsep\relax}{!}{%
\begin{tabular}{llllll}
\toprule
Epoch & Train Loss & Train Acc & Ratio & AUC & balance\_deviation \\
\midrule
1 & 0.3395 & 0.8588 & 6.513 & 0.952 & 0.400 \\
3 & 0.1041 & 0.9641 & 7.493 & 0.912 & 0.404 \\
**5** & **0.0733** & **0.9733** & **8.805** & **0.914** & **0.379** \\
10 & 0.0340 & 0.9908 & 8.509 & 0.888 & 0.374 \\
\bottomrule
\end{tabular}
}
\end{table*}

The ratio increases monotonically from epoch 1 (6.513) to epoch 5 (8.805) and remains at 8.509 at epoch 10. AUC is highest at epoch 1 (0.952) and declines slightly to 0.888 at epoch 10. The ratio at epoch 10 (8.509) substantially exceeds the 1.2x persistence threshold. The balance\_deviation metric decreases from 0.400 to 0.374 across epochs, reflecting a gradual shift in the $\tilde{g}$ distribution rather than any improvement relative to the criterion design flaw identified above.

The increase in ratio from epoch 1 to epoch 5 followed by a plateau at epoch 10 is consistent with majority group saturation: as majority samples converge under the spurious shortcut, their $\tilde{g} \rightarrow$ 0 and the ratio stabilizes. Whether the plateau would continue or reverse beyond epoch 10 was not measured.

\begin{figure}[htbp]
\centering
\includegraphics[width=\columnwidth]{figures/trajectory.png}
\caption{Temporal trajectory of proxy signal}
\label{fig:trajectory}
\end{figure}

\textit{Figure 3. Minority/majority $\tilde{g}$ ratio (left axis) and AUC (right axis) across epochs 1, 3, 5, and 10. The ratio increases from 6.5x to 8.8x across epochs 1--5 and remains at 8.5x at epoch 10.}

\subsubsection{5.4 Feature Norm Equalization}

The mechanistic interpretation of $\tilde{g}$ as a prediction-residual signal relies on $\|h(x_{i})\|$ being approximately uniform across groups. The h\_norm\_std\_ratio (ratio of standard deviation to mean of feature norms within each group) at T\_id = 5:

\begin{table}[htbp]
\centering
\resizebox{\columnwidth}{!}{%
\begin{tabular}{lll}
\toprule
Group & h\_norm mean & h\_norm\_std\_ratio \\
\midrule
G0 & 25.248 & $\approx$ 0.105 \\
G1 & 25.681 & $\approx$ 0.105 \\
G2 & 25.562 & $\approx$ 0.105 \\
G3 & 26.412 & $\approx$ 0.105 \\
\bottomrule
\end{tabular}
}
\end{table}

The h\_norm\_std\_ratio across all groups is approximately 0.105 at T\_id = 5 (ranging 0.097--0.118 across epochs 1--10), well below the 0.5 threshold. Feature norm means are within approximately 5\% across all four groups. This supports the interpretation that the $\tilde{g}$ signal reflects prediction residual differences rather than feature-scale differences between groups.

\begin{figure}[htbp]
\centering
\includegraphics[width=\columnwidth]{figures/feature_norms.png}
\caption{Feature norm distributions per group}
\label{fig:feature_norms}
\end{figure}

\textit{Figure 5. Feature norm $\|h(x_{i})\|$ distributions per group at T\_id = 5. Distributions are highly overlapping across groups (h\_norm\_std\_ratio $\approx$ 0.10), consistent with BatchNorm equalization.}

